\chapter{Vorgehensweise}
\label{cha:Vorgehensweise und Konzeption}

\textit{Planung, Auswahl der Komponenten und Designentscheidungen.}

% --- ANGUS ---
\section{Konstruktives Design und Hardware-Layout (Angus)}
\subsection{Konzept zur Stabilitätserhöhung}
\subsubsection{Positionierung der treibenden Räder und des Stützrades}
\subsubsection{Entwurf der Notausbremse (Hardwareseitig)}

\subsection{CAD-Entwicklung der Anbauteile}
\subsubsection{Planung der Roboterarme und Schultergelenke}
\subsubsection{Design des Roboterkopfes (Kamera- \& Displayaufnahme)}

% --- SIMON ---
\section{Softwarearchitektur der KI-Pipeline (Simon)}
\subsection{Konzept der Vollintegration}
\subsubsection{Datenfluss: Sprache LLM  Sprache}
\subsubsection{Planung der Gesichtsanimationen}

\subsection{Netzwerkinfrastruktur}
\subsubsection{Topologie: WLAN-Router, Repeater und Raspberry Pi}
\subsubsection{Konzept der Energieversorgung für Display und Pi}

% --- HAI NAM ---
\section{Elektromechanische Planung (Hai Nam)}
\subsection{Auslegung der Aktuatorik}
\subsubsection{Dimensionierung der Servos für Schulter und Klaue}
\subsubsection{Schaltplanentwurf: Screw Terminal und ESP32}

\subsection{Design- und Interaktionskonzept}
\subsubsection{Planung des Kabelbaums (Clean Design)}
\subsubsection{Logik-Entwurf der ``Speak-Through''-Funktion}

% --- FAIK ---
\section{Entwurf der Steuerungsalgorithmen (Faik)}
\subsection{Sicherheitslogik und Hinderniserkennung}
\subsubsection{Konzept der 3-Zonen-Ultraschallüberwachung}
\subsubsection{Definition der Fahr- und Bremsmodi}
\subsubsection{Anforderungsanalyse: Sicherheitslogik und Hinderniserkennung}

\subsection{Regelungskonzept für den Antrieb}
\subsubsection{Entwurf der Anfahrrampen (Exponentialfilter)}
\subsubsection{Geschwindigkeitsbegrenzung (Vorwärts vs. Rückwärts)}
\subsubsection{Anforderungsanalyse: Regelungstechnik für den Antrieb}